\documentclass[11pt,letterpaper]{article}
\usepackage[utf8]{inputenc}
\usepackage[margin=1in]{geometry}
\usepackage{amsmath}
\usepackage{amsfonts}
\usepackage{booktabs}
\usepackage{longtable}
\usepackage{hyperref}
\usepackage{listings}
\usepackage{xcolor}

\definecolor{codegreen}{rgb}{0,0.6,0}
\definecolor{codegray}{rgb}{0.5,0.5,0.5}
\definecolor{codepurple}{rgb}{0.58,0,0.82}
\definecolor{backcolour}{rgb}{0.95,0.95,0.92}
\definecolor{cautionred}{rgb}{0.55,0,0}
\definecolor{cautionbg}{rgb}{0.98,0.92,0.92}

\lstset{
    backgroundcolor=\color{backcolour},
    commentstyle=\color{codegreen},
    keywordstyle=\color{magenta},
    numberstyle=\tiny\color{codegray},
    stringstyle=\color{codepurple},
    basicstyle=\ttfamily\footnotesize,
    breakatwhitespace=false,
    breaklines=true,
    captionpos=b,
    keepspaces=true,
    numbers=left,
    numbersep=5pt,
    showspaces=false,
    showstringspaces=false,
    showtabs=false,
    tabsize=2
}

% Simple caution box. Deliberately does NOT capture its body as a macro
% argument (no \fcolorbox{}{}{...body...}) because several caution blocks
% contain lstlisting environments, and verbatim-like environments cannot
% be read correctly from inside a captured macro argument.
\newenvironment{caution}[1]{%
    \par\addvspace{8pt}%
    \noindent{\color{cautionred}\rule{\linewidth}{0.8pt}}\par\vspace{2pt}
    \noindent\textbf{\color{cautionred}Caution: #1}\par\vspace{2pt}
}{%
    \par\vspace{2pt}\noindent{\color{cautionred}\rule{\linewidth}{0.4pt}}%
    \par\addvspace{8pt}%
}

\title{Standard Operating Procedure (SOP): \\ Arctos Robotic Arm --- CANable Bring-Up, MKS Servo42D CAN Protocol, \\ and ROS~2 Joint Control (Joint B)}
\author{Assembled from project debugging logs, source-code inspection, \\ and empirical bench testing}
\date{August 2026}

\begin{document}

\maketitle
\tableofcontents
\newpage

\section{Introduction}

This guide walks through everything needed to take a fresh Ubuntu virtual machine and a CANable~2 USB-to-CAN adapter, and end up reliably talking to one joint (``Joint~B'') of an Arctos robotic arm over CAN bus, with a clear path toward integrating that control into the arm's ROS~2 driver.

It is written for someone who has never worked with a CAN bus or ROS~2 before. Two ideas make everything else in this guide make sense:

\subsection{CAN Bus, in one paragraph}
A Controller Area Network (CAN) bus is a two-wire (\texttt{CAN\_H} / \texttt{CAN\_L}) differential serial network. Every device on the bus --- in this project, every joint's motor controller --- listens to every message and only acts on the ones addressed to its own ID. Your computer does not talk CAN natively; the CANable~2 adapter converts USB traffic into real CAN electrical signals, and Linux's SocketCAN subsystem makes that adapter appear as an ordinary network interface (\texttt{can0}) that command-line tools and Python can use directly.

\subsection{Gearboxes and Closed-Loop Steppers, in one paragraph}
The motor at each joint does not drive the joint directly. It spins through a heavy reduction gearbox --- for Joint~B, a \textbf{67.82:1} ratio --- so the motor shaft must turn 67.82 times to rotate the physical joint once. The motor itself is a closed-loop NEMA~17 stepper: unlike a plain stepper, it has a magnetic encoder on the back of the shaft, so the MKS Servo42D controller always knows the true shaft position and can correct for missed steps. If the joint is obstructed and the shaft falls too far behind its commanded target, the controller trips a \emph{tracking-error protection stall} and locks the motor down; recovering from that requires power-cycling the controller.

\begin{caution}{Read this before touching real hardware}
Robotic arm joints, even small hobby ones, carry enough torque to pinch fingers or damage the linkage if commanded incorrectly. Keep hands clear of belts, pulleys, and linkages whenever the motor power supply is live, keep \texttt{candump} running in a separate terminal at all times, and know your disable/e-stop commands (Section~\ref{sec:safety}) before you send your first motion command.
\end{caution}

\newpage
\section{Hardware Overview}

\subsection{Joint / CAN ID Map}

\begin{table}[h]
\centering
\begin{tabular}{llll}
\toprule
\textbf{Joint} & \textbf{Physical Role} & \textbf{CAN ID (hex)} & \textbf{Bitrate} \\ \midrule
Joint A & Elbow & \texttt{0x04} & 500~kbit/s \\
Joint C & Wrist joint  & \texttt{0x06} & 500~kbit/s \\
Joint B & Wrist rotation & \texttt{0x05} & 500~kbit/s \\ \bottomrule
\end{tabular}
\caption{Arctos CAN node map. Only Joint~B (\texttt{0x05}) has been bench-tested in this guide; A and C are included for reference and use the same protocol.}
\end{table}

\subsection{Joint B Motor / Controller Specifications}

\begin{table}[h]
\centering
\begin{tabular}{ll}
\toprule
\textbf{Parameter} & \textbf{Value} \\ \midrule
Motor & NEMA 17, closed-loop stepper \\
Motor body length & 23~mm \\
Holding torque & 24~N\textperiodcentered cm \\
Rated current & 1.2~A \\
Shaft diameter & 5~mm \\
Motor step angle & 1.8\textdegree\ (200 full steps/revolution) \\
Controller & MKS Servo42D \\
Encoder resolution & 16{,}384 counts/revolution (14-bit) \\
Joint gear ratio & 67.82:1 \\
Joint direction & Inverted (\texttt{inverted: true} in config) \\ \bottomrule
\end{tabular}
\caption{Joint B hardware parameters, confirmed from \texttt{arctos\_controller.yaml} and the MKS datasheet-equivalent panel settings.}
\end{table}

\begin{caution}{Do not confuse motor steps with encoder counts}
The stepper's 200 full-steps/revolution is a physical stepping resolution. It is \emph{not} the same number as the encoder's 16{,}384 counts/revolution, which is what all CAN position math in this guide uses. Mixing the two up is a common source of a 80\texttimes\ scaling error.
\end{caution}

\subsection{CANable 2 Adapter}
\begin{itemize}
    \item USB VID:PID \texttt{16d0:117e}
    \item Contains its own CAN controller hardware --- \textbf{do not} add an external MCP2515 module in series with it.
    \item Only three wires are needed between the adapter and the MKS controller: \texttt{CAN\_H}, \texttt{CAN\_L}, and \texttt{GND}. Do \emph{not} connect the CANable's own power lines to the motor controller.
\end{itemize}

\subsection{Wiring Diagram}

\begin{lstlisting}[caption=Physical signal path]
Ubuntu Virtual Machine
        |
       USB
        |
    CANable 2
        |
 CAN_H / CAN_L / GND
        |
 MKS Servo42D Controller
        |
     Joint Motor
\end{lstlisting}

\subsection{Bus Termination Check (power off)}
With the system powered off, measure resistance between \texttt{CAN\_H} and \texttt{CAN\_L}:

\begin{table}[h]
\centering
\begin{tabular}{ll}
\toprule
\textbf{Reading} & \textbf{Meaning} \\ \midrule
$\approx 60\,\Omega$ & Correct --- two 120$\Omega$ terminators in parallel, wiring is likely fine \\
$\approx 120\,\Omega$ & Only one terminator present \\
Infinite / open & Wiring is disconnected \\
Very low & Possible short \\ \bottomrule
\end{tabular}
\end{table}

\newpage
\section{Bringing Up the CAN Interface on Ubuntu}
\label{sec:can-bringup}

Every time the CANable adapter is plugged in (including after a VirtualBox USB reattach), the interface must be rebuilt from scratch. This is the full sequence.

\subsection{Step 1: Confirm the USB device is visible}
\begin{lstlisting}[language=bash]
lsusb
# Expected line:
#   ID 16d0:117e CANable2
\end{lstlisting}

If using VirtualBox, also confirm the USB device is attached to the Ubuntu VM in the VirtualBox USB device menu, not just plugged into the host.

\subsection{Step 2: Find the serial device}
\begin{lstlisting}[language=bash]
ls -l /dev/ttyACM*
# Example: /dev/ttyACM0  (the number can change after every reconnect --
# always re-check this before starting CAN)
\end{lstlisting}

\subsection{Step 3: Load the CAN driver if needed}
\begin{lstlisting}[language=bash]
lsmod | grep gs_usb
# If nothing prints:
sudo modprobe gs_usb
sudo dmesg | tail -30
# Look for: "registered new interface driver gs_usb"
\end{lstlisting}

\subsection{Step 4: Clean up any old session}
\begin{lstlisting}[language=bash]
sudo pkill slcand
sudo ip link delete can0 2>/dev/null
\end{lstlisting}

\subsection{Step 5: Bind the adapter to a CAN interface}

\begin{caution}{The slcand bitrate flag is a common mistake}
\texttt{slcand}'s \texttt{-s} flag selects the bitrate by a fixed code, \emph{not} the raw number: \texttt{-s0}=10k, \texttt{-s1}=20k, \texttt{-s2}=50k, \texttt{-s3}=100k, \texttt{-s4}=125k, \texttt{-s5}=250k, \textbf{\texttt{-s6}=500k}, \texttt{-s7}=800k, \texttt{-s8}=1M. Joint~B is configured for \textbf{500~kbit/s}, so the correct flag is \texttt{-s6}. Earlier debugging logs for this project inconsistently used \texttt{-s8} (which actually requests 1~Mbit/s) in some sessions --- if a bitrate mismatch happens, \texttt{candump} will show your own outgoing frames but the motor will never respond (see Section~\ref{sec:troubleshooting}). Always double-check this flag, or set the bitrate explicitly in Step 6 below to be safe either way.
\end{caution}

\begin{lstlisting}[language=bash]
sudo slcand -o -c -s6 /dev/ttyACM0 can0
\end{lstlisting}

\subsection{Step 6: Bring the link up at the correct bitrate}
Setting the bitrate explicitly on the \texttt{ip link} command as well removes any ambiguity from the \texttt{slcand} flag:
\begin{lstlisting}[language=bash]
sudo ip link set can0 up type can bitrate 500000
\end{lstlisting}

\subsection{Step 7: Verify}
\begin{lstlisting}[language=bash]
ip -details -statistics link show can0
\end{lstlisting}
Look for:
\begin{lstlisting}
can0: <NOARP,UP,LOWER_UP> ...
bitrate 500000
can state ERROR-ACTIVE
\end{lstlisting}
\texttt{RX packets} should increase whenever the motor responds to something; \texttt{TX packets} increase whenever you send. Bus errors climbing, or a \texttt{bus-off} state, both point back to wiring or bitrate.

\subsection{Shutdown procedure}
Before unplugging the CANable adapter:
\begin{lstlisting}[language=bash]
# 1. Stop any running Python control script: Ctrl+C
# 2. Bring the interface down
sudo ip link set can0 down
# 3. Stop slcand
sudo pkill slcand
# 4. Then physically unplug the adapter
\end{lstlisting}

\newpage
\section{MKS Servo42D Panel Configuration}

On the controller itself (Joint~B), confirm these settings before attempting CAN communication:

\begin{table}[h]
\centering
\begin{tabular}{ll}
\toprule
\textbf{Setting} & \textbf{Value} \\ \midrule
CAN ID & \texttt{05} \\
CAN Rate & \texttt{500} (kbit/s) \\
CAN Response (CANRSP) & \texttt{ENABLE} \\
CAN CRC & \texttt{SUM-8} \\
Working Mode & \texttt{CR\_CLOSE} (see caution below) \\ \bottomrule
\end{tabular}
\end{table}

Save settings and power-cycle the controller after any change. The CAN ID in particular must actually be committed, not just visible in the menu.

\begin{caution}{Working-mode mismatch is a leading suspect for ``command acknowledged, motor doesn't move''}
The Arctos ROS~2 driver source defines six working modes and defaults to \texttt{working\_mode = 2}, i.e.\ \texttt{CR\_vFOC}:
\begin{lstlisting}[language=C++]
enum class MotorMode {
    CR_OPEN  = 0,
    CR_CLOSE = 1,
    CR_vFOC  = 2,
    SR_OPEN  = 3,
    SR_CLOSE = 4,
    SR_vFOC  = 5
};
\end{lstlisting}
If the physical panel is set to \texttt{CR\_CLOSE} while a motion command was generated assuming \texttt{CR\_vFOC} (or vice versa), the controller can accept and checksum-acknowledge a motion frame without actually rotating the shaft --- which matches symptoms observed during testing (motor made a sound but did not visibly rotate). \textbf{Confirm the panel's actual working mode and the value used when constructing \texttt{0x82} (\texttt{SET\_WORKING\_MODE}) frames before trusting any ``no motion'' or ``no response'' result as a wiring/CAN problem.}
\end{caution}

\newpage
\section{The MKS CAN Protocol as Used by Arctos}

\subsection{Frame Structure and Checksum}

Every command frame sent to Joint~B follows the same pattern:

\begin{center}
\texttt{[command byte] [data bytes...] [checksum byte]}
\end{center}

The checksum has been independently verified (matched by hand against multiple successful transmissions) as a simple 8-bit sum:

\begin{lstlisting}[language=python]
def checksum(motor_id, data_bytes):
    return (motor_id + sum(data_bytes)) & 0xFF
\end{lstlisting}

\begin{table}[h]
\centering
\begin{tabular}{lll}
\toprule
\textbf{Frame (no checksum)} & \textbf{Arithmetic} & \textbf{Checksum} \\ \midrule
\texttt{31} & \texttt{0x31+0x05} & \texttt{36} \\
\texttt{3A} & \texttt{0x3A+0x05} & \texttt{3F} \\
\texttt{F1} & \texttt{0xF1+0x05} & \texttt{F6} \\
\texttt{F3 01} & \texttt{0xF3+0x01+0x05} & \texttt{F9} \\
\texttt{F5 00 64 14 00 0C E4} & sum $+$ \texttt{0x05} & \texttt{62} \\ \bottomrule
\end{tabular}
\caption{Worked checksum examples for CAN ID \texttt{0x05} (Joint B). All verified by hand and against actual bus traffic.}
\end{table}

\begin{caution}{Never hand-invent a checksum}
Manually sent \texttt{cansend} commands with an incorrect or omitted checksum byte were the direct cause of several ``the motor ignored my command'' failures earlier in this project. Always compute the checksum from the formula above, never guess it.
\end{caution}

\subsection{Command Reference}

\begin{longtable}{lll}
\toprule
\textbf{Name} & \textbf{Byte} & \textbf{Notes} \\ \midrule
\endhead
\texttt{QUERY\_MOTOR} & \texttt{0xF1} & Confirms the controller is alive. Verified response: \texttt{F1 01 F7}. \\
\texttt{READ\_ENCODER} (carry form) & \texttt{0x30} & See Section~\ref{sec:encoder}. \\
\texttt{READ\_ENCODER} (alt.\ form) & \texttt{0x31} & Different response layout observed --- unverified, see Section~\ref{sec:encoder}. \\
\texttt{READ\_VELOCITY} & \texttt{0x32} & Response layout not yet confirmed; log raw bytes only. \\
\texttt{READ\_ERROR} & \texttt{0x39} & Response layout not yet confirmed. \\
\texttt{READ\_ENABLE\_STATE} & \texttt{0x3A} & Response byte-to-state mapping not yet confirmed. \\
\texttt{READ\_SHAFT\_PROTECTION\_STATE} & \texttt{0x3E} & Indicates a tracking-error stall. \\
\texttt{SET\_ENABLE\_SETTINGS} & \texttt{0x85} & Configuration, not a runtime enable --- do not confuse with \texttt{0xF3}. \\
\texttt{SET\_WORKING\_MODE} & \texttt{0x82} & See working-mode caution above. \\
\texttt{ENABLE\_MOTOR} & \texttt{0xF3} & Payload \texttt{0x01}=enable, \texttt{0x00}=disable. Confirmed working both directions. \\
\texttt{RELATIVE\_POSITION} & \texttt{0xF4} & See Section~\ref{sec:motion}. \\
\texttt{ABSOLUTE\_POSITION} & \texttt{0xF5} & See Section~\ref{sec:motion}. \\
\texttt{SPEED\_CONTROL} & \texttt{0xF6} & Not exercised in testing so far. \\
\texttt{EMERGENCY\_STOP} & \texttt{0xF7} & See Section~\ref{sec:safety}. \\
\texttt{POSITION\_CONTROL} & \texttt{0xFD} & Not exercised in testing so far. \\
\texttt{ABSOLUTE\_POSITION\_PULSE} & \texttt{0xFE} & Not exercised in testing so far. \\
\bottomrule
\caption{Command bytes confirmed from \texttt{motor\_types.hpp}. ``Not yet confirmed'' items should be logged with raw \texttt{candump} output and cross-checked against source before writing parsing code.}
\end{longtable}

\subsection{Enable / Disable / Query}

\begin{lstlisting}[language=bash, caption={Enable, disable, and query -- all checksum-verified}]
# Enable
cansend can0 005#F301F9

# Disable
cansend can0 005#F300F8

# Query (liveness check)
cansend can0 005#F1F6
# Expected response: 005#F1 01 F7
\end{lstlisting}

\subsection{Reading the Encoder}
\label{sec:encoder}

The current ROS~2 source explicitly defines \texttt{READ\_ENCODER = 0x31} in
\texttt{motor\_types.hpp}.  The receive path in \texttt{motor\_driver.cpp} sends
the six bytes following the \texttt{0x31} command byte to
\texttt{CANProtocol::decodeInt48()}.  Therefore, for the current driver,
\textbf{\texttt{0x31} is the encoder-position command used for position math}.

\begin{lstlisting}[language=bash, caption={Observed Joint B encoder traffic}]
cansend can0 005#3136
# Observed response:
# can0 005 [8] 31 FF FF FF FF FF F7 28

cansend can0 005#3136
# Another observed response:
# can0 005 [8] 31 00 00 00 00 00 01 37
\end{lstlisting}

The current C++ decoder requires exactly six payload bytes after the command:
bytes 1--6 of the CAN frame are combined as one signed 48-bit big-endian
integer.  The decoder then converts that value to degrees using

\begin{lstlisting}[language=C++]
int64_t addition_value =
    (static_cast<int64_t>(data[0]) << 40) |
    (static_cast<int64_t>(data[1]) << 32) |
    (static_cast<int64_t>(data[2]) << 24) |
    (static_cast<int64_t>(data[3]) << 16) |
    (static_cast<int64_t>(data[4]) << 8)  |
    static_cast<int64_t>(data[5]);

if (addition_value & 0x800000000000) {
    addition_value |= 0xFFFF000000000000;
}

double motor_angle_deg =
    (addition_value * MotorConstants::DEGREES_PER_REVOLUTION)
    / MotorConstants::ENCODER_STEPS;
\end{lstlisting}

The driver's unit tests provide the intended scale: a six-byte value of
\texttt{00 00 01 00 00 00} must decode to $360^\circ$, and
\texttt{00 00 00 00 20 00} must decode to $180^\circ$.  Thus the source
currently treats \texttt{0x4000 = 16384} encoder counts as one motor-shaft
revolution.

\begin{caution}{Do not mix the three observed read commands}
The current source defines \texttt{READ\_ENCODER = 0x31},
\texttt{READ\_PULSES = 0x33}, and \texttt{READ\_RAW\_ENCODER = 0x35} as
different commands.  Captured responses include:

\begin{lstlisting}[language=bash]
can0 005 [8] 31 00 00 00 00 00 01 37
can0 005 [6] 33 00 00 00 23 5B
can0 005 [8] 35 00 00 00 00 0B 62 A7
\end{lstlisting}

Only the \texttt{0x31} response is passed through
\texttt{decodeInt48()} by the current ROS~2 driver.  The exact numerical
meaning of the \texttt{0x33} and \texttt{0x35} payloads has not been
established from the source examined so far, so they should not be substituted
into the \texttt{0x31} position formula.
\end{caution}

\subsection{Motion Commands}
\label{sec:motion}

The Arctos driver's absolute-position frame construction, confirmed directly from \texttt{motor\_driver.cpp}:

\begin{lstlisting}[language=C++]
double motor_position = position * joint.gear_ratio;
double motor_position_deg = motor_position * MotorConstants::RAD_TO_DEG;
int32_t encoder_counts = static_cast<int32_t>(
    (motor_position_deg * MotorConstants::ENCODER_STEPS)
    / MotorConstants::DEGREES_PER_REVOLUTION
);

std::vector<uint8_t> data = {
    CANCommands::ABSOLUTE_POSITION,                         // 0xF5
    static_cast<uint8_t>((speed >> 8) & 0xFF),
    static_cast<uint8_t>(speed & 0xFF),
    acc_value,
    static_cast<uint8_t>((encoder_counts >> 16) & 0xFF),
    static_cast<uint8_t>((encoder_counts >> 8) & 0xFF),
    static_cast<uint8_t>(encoder_counts & 0xFF)
};
\end{lstlisting}

Frame layout: \texttt{[0xF5][speed\_hi][speed\_lo][accel][pos\_hi][pos\_mid][pos\_lo][checksum]}. The 3-byte position field is wide enough (24 bits, up to $\approx$16.7~million) to carry a full multi-turn absolute count directly -- it is \emph{not} limited to one revolution.

\begin{caution}{Absolute-vs-relative behavior of \texttt{0xF5} is not fully confirmed}
Two bench tests sent structurally valid, correctly checksummed \texttt{0xF5} frames and produced results that do not cleanly agree with a single interpretation: one test's resulting motion was closer to what a \emph{relative delta} would produce, the next test's was not. Until this is resolved by reading \texttt{can\_protocol.cpp}'s \texttt{sendFrame()} implementation and comparing it against controlled single-variable tests, treat any \texttt{0xF5} target as \textbf{unverified}, and always keep the commanded magnitude small enough that either interpretation is safe (Section~\ref{sec:safe-test}).
\end{caution}

\newpage
\section{From Raw Encoder Counts to Joint Degrees}

Combining the confirmed encoder resolution and gear ratio:

\[
\text{counts per Joint-B degree} = \frac{\text{gear\_ratio} \times \text{ENCODER\_STEPS}}{360} = \frac{67.82 \times 16384}{360} \approx 3086.56
\]

\begin{table}[h]
\centering
\begin{tabular}{ll}
\toprule
\textbf{Joint B motion} & \textbf{Raw encoder counts (approx.)} \\ \midrule
$0.01^\circ$ & 31 \\
$0.1^\circ$ & 309 \\
$1^\circ$ & 3087 \\
$360^\circ$ (one full joint revolution) & 1{,}111{,}163 \\ \bottomrule
\end{tabular}
\caption{Worked conversions using the verified 67.82:1 gear ratio and 16{,}384-count encoder.}
\end{table}

\begin{caution}{Two other conversion numbers exist in this project's history -- do not use them}
An early calibration pass, done before the gear ratio was discovered in \texttt{arctos\_controller.yaml}, assumed the motor's raw shaft rotation \emph{was} the joint rotation and produced a figure of ``$\approx$3300 counts per 90$^\circ$.'' That number is real data, but it describes the \emph{motor shaft}, not the joint -- once the 67.82:1 reduction is applied, that same 3300-count motion corresponds to only $\approx 1.07^\circ$ of actual Joint~B travel. Separately, a since-superseded Python calibration note referenced
\texttt{COUNTS\_PER\_REV~=~16050}.  That value is not used by the current
C++ driver.  The current source explicitly defines
\texttt{ENCODER\_STEPS~=~16384} and the unit tests are written around that
scale.  For calculations that reproduce the current ROS~2 driver's command
math, use the source-defined 16384 value.  The 16050 value should be retained
as historical empirical data, not silently substituted into the C++ formula.
\end{caution}

\begin{caution}{A conflicting ``$-813$ pulses/degree'' figure was also found -- do not use it}
A separate driver draft (\texttt{mks\_driver.py}) defined
\texttt{pulses\_per\_degree~=~-813.0} for Joint~B and used it directly as the
value written into the CAN frame's position field.  Its origin has not yet
been reconciled with the current ROS~2 source:
\begin{lstlisting}[language=python]
target_pulses = int(degrees * profile["pulses_per_degree"])
estimated_final_encoder = current_encoder + int(target_pulses * -0.947)
MIN_LIMIT, MAX_LIMIT = -73000, 73000
if estimated_final_encoder < MIN_LIMIT or estimated_final_encoder > MAX_LIMIT:
    print("Movement Blocked...")
    return
pos = target_pulses & 0xFFFFFF   # <- this is what is actually transmitted
\end{lstlisting}
This does not match the source-verified $\approx 3086.56$ counts/degree figure -- it is off by roughly a factor of 3.8. More importantly, the safety-limit \emph{estimate} (\texttt{target\_pulses~*~-0.947}) uses a completely different multiplier than the value that is actually transmitted on the bus (\texttt{target\_pulses} directly). Because the two don't agree, the ``Safety Workspace Cage'' cannot reliably predict where a command will actually send the joint, and so cannot be trusted to block an unsafe one. \textbf{Do not use this conversion constant or this safety-limit code until it is reconciled against the verified gear-ratio math above.}
\end{caution}

\newpage
\section{Safety Commands}
\label{sec:safety}

\begin{table}[h]
\centering
\begin{tabular}{lll}
\toprule
\textbf{Action} & \textbf{Command} & \textbf{Notes} \\ \midrule
Disable motor & \texttt{cansend can0 005\#F300F8} & Checksum-verified; use this form. \\
Enable motor & \texttt{cansend can0 005\#F301F9} & Checksum-verified. \\
Emergency stop & \texttt{cansend can0 005\#F7FC} & \textbf{Verify before relying on it -- see caution below.} \\ \bottomrule
\end{tabular}
\end{table}

\begin{caution}{Confirm the emergency-stop frame before you need it in an emergency}
Notes from this project show the emergency-stop command sent as a bare \texttt{F7} with no checksum byte at all, which is inconsistent with every other verified command on this bus (all of which required a checksum, \texttt{0xF7+0x05=0xFC} by the same formula used everywhere else). Do not assume either form works under load. With the joint mechanically unloaded and clear of obstructions, deliberately trigger both forms and confirm with \texttt{candump} that the motor actually loses holding torque before you depend on either one during a real motion test.
\end{caution}

\newpage
\section{A Safe First Motion Test}
\label{sec:safe-test}

Given the open questions in Sections~\ref{sec:motion} and above, the right first motion test is one that is safe \emph{regardless of which interpretation turns out to be correct}. A magnitude of $\approx 300$ raw counts ($\approx 0.1^\circ$ of Joint~B) satisfies this: if the firmware treats the field as an absolute target, the worst case is a small jump near zero; if it treats it as a relative delta, it is simply a $0.1^\circ$ nudge.

\subsection{Procedure}
\begin{enumerate}
    \item Bring up \texttt{can0} per Section~\ref{sec:can-bringup}.
    \item Open a second terminal and leave \texttt{candump -tz can0} running.
    \item Confirm the joint is mechanically clear.
    \item Run the script below. It reads the true absolute position (carry-aware), pauses for confirmation, sends one small move, re-reads the position, and reports the actual counts and degrees moved -- which also tells you empirically whether \texttt{0xF5} is behaving as absolute or relative.
    \item The motor is disabled in a \texttt{finally} block regardless of how the script exits.
\end{enumerate}

\subsection{Script: \texttt{joint\_b\_micro\_move\_test.py}}
\begin{lstlisting}[language=python]
#!/usr/bin/env python3
"""
Safe, small (~0.1 deg) test move for Arctos Joint B (MKS Servo42D, CAN ID 0x05).
Reads the FULL absolute encoder (carry * 16384 + value) before and after,
uses a magnitude small enough to be safe under either an absolute or
relative interpretation of the 0xF5 position field, and always disables
the motor before exiting.
"""
import can
import time

CAN_ID = 0x05
CHANNEL = "can0"
BUSTYPE = "socketcan"

TEST_DELTA_COUNTS = 300      # ~0.1 degree of Joint B (67.82:1 gear, 16384 cpr)
SPEED = 20                   # low speed for safety
ACCEL = 5                    # low acceleration for safety
COUNTS_PER_DEGREE = 67.82 * 16384 / 360.0   # ~3086.56


def checksum(motor_id, data_bytes):
    return (motor_id + sum(data_bytes)) & 0xFF


def send_frame(bus, motor_id, data_bytes):
    crc = checksum(motor_id, data_bytes)
    frame = data_bytes + [crc]
    msg = can.Message(arbitration_id=motor_id, data=frame, is_extended_id=False)
    bus.send(msg)
    return frame


def read_encoder(bus, motor_id, timeout=1.0):
    send_frame(bus, motor_id, [0x30])
    end = time.time() + timeout
    while time.time() < end:
        resp = bus.recv(timeout=end - time.time())
        if resp is None:
            break
        if resp.arbitration_id == motor_id and len(resp.data) >= 7 and resp.data[0] == 0x30:
            carry = int.from_bytes(resp.data[1:5], byteorder="big", signed=True)
            value = int.from_bytes(resp.data[5:7], byteorder="big", signed=False)
            return carry, value, carry * 16384 + value
    raise TimeoutError("No encoder response received")


def enable_motor(bus, motor_id):
    send_frame(bus, motor_id, [0xF3, 0x01]); time.sleep(0.2)


def disable_motor(bus, motor_id):
    send_frame(bus, motor_id, [0xF3, 0x00]); time.sleep(0.2)


def move_small(bus, motor_id, delta_counts, speed, accel):
    speed_hi, speed_lo = (speed >> 8) & 0xFF, speed & 0xFF
    pos = delta_counts & 0xFFFFFF
    pos_hi, pos_mid, pos_lo = (pos >> 16) & 0xFF, (pos >> 8) & 0xFF, pos & 0xFF
    data = [0xF5, speed_hi, speed_lo, accel, pos_hi, pos_mid, pos_lo]
    return send_frame(bus, motor_id, data)


def main():
    bus = can.interface.Bus(channel=CHANNEL, bustype=BUSTYPE)
    try:
        carry0, value0, abs0 = read_encoder(bus, CAN_ID)
        print(f"Start: carry={carry0} value={value0} absolute={abs0}")

        input("Enable motor and send a ~0.1 degree test move? "
              "Press Enter to continue, Ctrl+C to abort...")

        enable_motor(bus, CAN_ID)
        frame = move_small(bus, CAN_ID, TEST_DELTA_COUNTS, SPEED, ACCEL)
        print(f"Sent frame: {[hex(b) for b in frame]}")
        time.sleep(2.0)

        carry1, value1, abs1 = read_encoder(bus, CAN_ID)
        moved = abs1 - abs0
        print(f"End:   carry={carry1} value={value1} absolute={abs1}")
        print(f"Observed change: {moved} raw counts "
              f"({moved / COUNTS_PER_DEGREE:.4f} deg of Joint B)")
        print("moved ~= 300      -> field behaves as a RELATIVE delta")
        print("moved ~= 300-abs0 -> field behaves as an ABSOLUTE target")
    finally:
        print("Disabling motor.")
        disable_motor(bus, CAN_ID)
        bus.shutdown()


if __name__ == "__main__":
    main()
\end{lstlisting}

\newpage
\section{ROS 2 Integration (Project Layout and Status)}

The confirmed working CAN-level control (Sections~\ref{sec:can-bringup}--\ref{sec:safe-test}) is the foundation for the ROS~2 driver package. As of this guide, \textbf{full closed-loop joint control through ROS~2 has not yet been demonstrated} -- only direct, script-level CAN control has. The pieces below are the confirmed structure to build on.

\subsection{Current source-verified command and conversion definitions}
The current \texttt{motor\_types.hpp} defines the following relevant values:

\begin{lstlisting}[language=C++]
READ_ENCODER          = 0x31
READ_PULSES           = 0x33
READ_RAW_ENCODER      = 0x35
SET_WORKING_MODE      = 0x82
SET_SUBDIVISIONS      = 0x84
SET_ENABLE_SETTINGS   = 0x85

ENABLE_MOTOR          = 0xF3
RELATIVE_POSITION     = 0xF4
ABSOLUTE_POSITION     = 0xF5
POSITION_CONTROL      = 0xFD
ABSOLUTE_POSITION_PULSE = 0xFE

MotorParameters::working_mode  = 2
MotorParameters::subdivisions  = 16

MotorConstants::ENCODER_STEPS = 16384.0
\end{lstlisting}

The mode enumeration identifies mode~2 as \texttt{CR\_vFOC}, while the default
\texttt{MotorParameters::working\_mode} is also 2.  This documents what the
current code assumes; it does \emph{not} establish that the physical motor
panel is currently configured to that mode.

\subsection{Package layout}
\begin{lstlisting}[language=bash]
~/arctos_ws/src/ros2_arctos/arctos_motor_driver/
  include/arctos_motor_driver/motor_types.hpp   # command byte + mode enums
  src/motor_driver.cpp                          # frame construction, gear-ratio math
  src/can_protocol.cpp                          # sendFrame() / checksum
~/arctos_ws/src/ros2_arctos/arctos_description/config/arctos_controller.yaml
  # per-joint gear_ratio, inverted, zero_position, home_position, opposite_limit
\end{lstlisting}

\subsection{Joint B's configuration block}
\begin{lstlisting}
B_joint:
    motor_id: 5
    requires_homing: false
    inverted: true
    gear_ratio: 67.82
    hardware_type: MKS_42D
    zero_position: -0.43627526433400465    # radians
    home_position: 0.22467435360166688     # radians
    opposite_limit: -0.4228851636789132    # radians
\end{lstlisting}
These are joint-space parameters in radians, not raw encoder counts -- do not substitute a raw encoder reading for one of these without going through the gear-ratio conversion in Section 5 of this guide.

\subsection{Current CAN checksum implementation}
The current \texttt{CANProtocol::sendFrame()} implementation sets the DLC to
the command-data length plus one checksum byte, copies the command data into
the CAN frame, and computes the checksum as the low eight bits of the sum of
the motor ID and every command-data byte:

\begin{lstlisting}[language=C++]
uint16_t crc = motor_id;
for (size_t i = 0; i < data.size(); i++) {
    crc += data[i];
}
crc &= 0xFF;
msg->data[data.size()] = static_cast<uint8_t>(crc);
\end{lstlisting}

This matches the checksum convention used elsewhere in this guide for the
frames examined so far.

\subsection{Building and running}
\begin{lstlisting}[language=bash]
cd ~/arctos_ws
colcon build --packages-select arctos_motor_driver
source ~/arctos_ws/install/setup.bash

ros2 node list
ros2 topic list
ros2 topic echo /topic_name
ros2 service list
ros2 action list
\end{lstlisting}

\subsection{Recommended next steps toward full ROS 2 control}
\begin{enumerate}
    \item Resolve the \texttt{0xF5} absolute-vs-relative question using controlled,
          single-variable motion tests.
    \item Determine the exact semantics of the observed \texttt{0x33} and
          \texttt{0x35} responses before using either as a substitute for the
          \texttt{0x31} position response.
    \item Trace the origin of the historical \texttt{-813 pulses/degree} value
          in Git history / older scripts.  Do not use it merely because it
          appears in an older driver draft.
    \item Confirm the physical motor panel's working mode against the source's
          default \texttt{working\_mode=2}; the source definition alone does
          not prove the panel setting.
    \item Only once single-joint CAN control is deterministic, wire it into the
          ROS~2 controller's command topic / action interface.
\end{enumerate}

\newpage
\section{Troubleshooting}
\label{sec:troubleshooting}

\begin{longtable}{p{0.32\textwidth}p{0.6\textwidth}}
\toprule
\textbf{Symptom} & \textbf{Likely cause / fix} \\ \midrule
\endhead
CANable disappears after unplugging & \texttt{/dev/ttyACM\#} number changed. Re-run \texttt{lsusb} and \texttt{ls -l /dev/ttyACM*}, use the new number in \texttt{slcand}. \\
\texttt{can0} does not exist & Old \texttt{slcand} session still holding it. \texttt{sudo pkill slcand; sudo ip link delete can0} then recreate. \\
\texttt{candump} only shows your own outgoing frames & Adapter is transmitting but nothing replies. Check bitrate (Caution in Section~\ref{sec:can-bringup}), CAN ID, CANRSP=ENABLE, and \texttt{CAN\_H}/\texttt{CAN\_L} wiring. \\
MKS display is on, no CAN response at all & Confirm CAN ID, CAN rate, CANRSP, and working mode on the panel match what you are sending. \\
Command is checksummed correctly and acknowledged, but motor does not move & Suspect working-mode mismatch (Section~3) or an incorrect \texttt{0xF5}/\texttt{0xF4} interpretation (Section~\ref{sec:motion}) before assuming a wiring fault. \\
Motor gets hot during repeated testing & Expected for a continuously-energized stepper -- disable (\texttt{F300F8}) between tests and let it cool; do not leave it enabled unnecessarily. \\
\bottomrule
\caption{Common failure modes observed during Joint B bring-up.}
\end{longtable}

\newpage
\section{Quick Reference}

\begin{lstlisting}[language=bash, caption=Copy/paste reference for a normal session]
# --- Bring up CAN ---
lsusb
ls -l /dev/ttyACM*
sudo pkill slcand
sudo ip link delete can0 2>/dev/null
sudo slcand -o -c -s6 /dev/ttyACM0 can0
sudo ip link set can0 up type can bitrate 500000
ip -details -statistics link show can0

# --- Monitor ---
candump -tz can0

# --- Query / read (all checksum-verified) ---
cansend can0 005#F1F6      # query motor        -> F1 01 F7
cansend can0 005#3035      # read encoder        -> 30 <carry:4><value:2><crc>

# --- Enable / disable / stop ---
cansend can0 005#F301F9    # enable
cansend can0 005#F300F8    # disable
cansend can0 005#F7FC      # emergency stop (VERIFY before relying on it)

# --- Shutdown ---
sudo ip link set can0 down
sudo pkill slcand
\end{lstlisting}

\end{document}